\section{Hoja de ruta de implementación en ingeniería: de los paradigmas de la clase al experimento reproducible}

\subsection{Primero se define el «objeto optimizable», después se optimiza el modelo}
Lo más útil de trasladar esta clase al desarrollo cotidiano es lo siguiente: \textbf{primero hay que definir la tarea como un objeto ejecutable, verificable e iterable, y solo después hablar de optimizar el modelo}. Muchos equipos fracasan no porque el modelo no sea lo bastante potente, sino porque la definición de la tarea está incompleta; por ejemplo, solo cuentan con un objetivo en lenguaje natural sin un verifier explícito, o solo con la puntuación final sin diagnóstico a nivel de trayectoria.

\begin{importantbox}{La unidad experimental mínima de un Code Agent de seguridad}
Un experimento reproducible debe incluir al menos seis componentes:
\begin{enumerate}
    \item especificación de la tarea: entrada, salida, restricciones, tiempo límite;
    \item entorno de ejecución: contenedor reconstruible y versiones de las dependencias;
    \item interfaz de herramientas: exploración de código, ejecución, depuración, generación de entradas;
    \item verificador: conjunto de pruebas o sanitizer/oracle;
    \item registro de trayectoria: thought/tool/observation en cada ronda;
    \item script de evaluación: ejecución por lotes y agregación estadística.
\end{enumerate}
\end{importantbox}

Si faltan los puntos 5 y 6, el sistema suele quedar en la situación de ``ver solo el resultado, sin ver el proceso'', lo que dificulta reproducir los problemas. Aquí es donde realmente aterriza lo que Sutton dice con ``always look at the data'': el equipo debe poder tratar las trayectorias fallidas como un activo de datos de primer orden y analizarlas de manera continua.

\begin{knowledgebox}{Estrategia de puente entre la clase y la ingeniería}
La estructura en tres partes del curso puede trasladarse a un flujo de desarrollo de tres capas:
\begin{itemize}
    \item capa base: evaluación de tareas de código (MBPP/HumanEval/SWE-Bench) para establecer las capacidades fundamentales;
    \item capa de sistema: harness, herramientas y flujo de control para mejorar la estabilidad en tareas complejas;
    \item capa de aplicación: verificación dinámica y cadena de evidencia en tareas de seguridad para comprobar el valor real.
\end{itemize}
\end{knowledgebox}

\subsection{Ritmo de iteración: descomponer los cuellos de botella del sistema por tipo de fallo}
En clase se mencionó varias veces la ``optimización combinatoria'', que en ingeniería puede llevarse un paso más allá mediante la clasificación de fallos. En cuanto las causas de los fallos se clasifican de forma estable, las prioridades de ajuste quedan mucho más claras, en lugar de reemplazar el modelo a ciegas o apilar prompts.

\begin{table}[H]
\centering
\begin{tabular}{p{3.2cm}p{4.6cm}p{4.6cm}}
\toprule
\textbf{Tipo de fallo} & \textbf{Síntoma típico} & \textbf{Dirección prioritaria de corrección} \\
\midrule
Fallo de recuperación & Explora repetidamente archivos irrelevantes, la localización de la ruta se desvía & Mejorar la búsqueda a nivel de símbolo, incorporar navegación por referencias cruzadas \\
Fallo de ejecución & Formato incorrecto en la invocación de comandos, dependencias de entorno faltantes & Restringir los parámetros de las herramientas, fijar una línea base del contenedor \\
Fallo de verificación & Las pruebas pasan pero hay error semántico, o se juzga mal la causa del fallo & Reforzar el verifier, añadir repetición de contraejemplos \\
Fallo de estrategia & En trayectorias largas se autorrefuerzan hipótesis erróneas & Reiniciar el flujo de control, ciclos por etapas, restricciones de máquina de estados \\
Fallo de costo & La tasa de éxito mejora, pero el consumo de tokens o la latencia son inaceptables & Gestión del presupuesto, estrategias de parada temprana, enrutamiento de modelos por niveles \\
\bottomrule
\end{tabular}
\caption{Mapeo de los modos de fallo del Agente a direcciones de mejora del sistema que se puedan poner en práctica}
\end{table}

\begin{warningbox}{Revisar solo «casos de éxito» induce al equipo a error}
Si en las reuniones de revisión solo se muestra la best trajectory, el sistema terminará optimizándose para ser ``ocasionalmente impresionante''. Para mejorar la fiabilidad en producción es necesario revisar de manera sistemática las muestras de fallo de cola larga, en especial las trayectorias que ``parecen a punto de tener éxito pero terminan colapsando''.
\end{warningbox}

\subsection{Gobernanza del flujo de control: el equilibrio entre eficiencia de exploración y capacidad de auditoría}
Según el espectro continuo que presenta el curso, el flujo de control no consiste en tomar partido por un extremo, sino en configurarlo según el escenario. La fase de exploración puede inclinarse hacia lo dinámico, y la fase de convergencia hacia lo restringido. En las tareas de seguridad, además, se deben cumplir requisitos de auditoría: que cada acción clave se pueda rastrear y que cada conclusión se pueda reproducir.

\begin{importantbox}{Flujo de control de doble modo recomendado}
\begin{itemize}
    \item modo de exploración: muestreo alto, múltiples rutas, permisos amplios de herramientas, para descubrir posibles rutas de vulnerabilidad;
    \item modo de verificación: muestreo bajo, restricciones estrictas, puntos de observación fijos, para reproducir y consolidar evidencia.
\end{itemize}
Ambos modos comparten el mismo objetivo de la tarea, pero difieren en la estrategia de recursos y en los requisitos de interpretabilidad.
\end{importantbox}

\begin{knowledgebox}{Reglas de gestión de presupuesto que se pueden poner en práctica}
El presupuesto en tiempo de prueba puede dividirse en tres capas:
\begin{itemize}
    \item presupuesto por trayectoria: limita el número máximo de pasos y de ejecuciones de cada trayectoria;
    \item presupuesto por tarea: limita el número total de trayectorias y la duración total de cada ejemplo;
    \item presupuesto por lote: limita el total de horas de GPU/CPU de toda la ronda de evaluación.
\end{itemize}
Cuando alguna capa alcanza su umbral, el sistema debe generar una razón estructurada de ``por qué se detuvo'', para evitar que el tiempo de espera agotado sea una caja negra.
\end{knowledgebox}

\subsection{Restricciones adicionales para la implementación en seguridad: divulgación responsable y riesgo de doble uso}
El caso Big Sleep demuestra que ``descubrir una vulnerabilidad'' ya es viable en la práctica, por lo que, al ponerlo en producción, no basta con hablar de métricas técnicas. Un Agente de seguridad necesita límites de responsabilidad integrados, entre ellos un proceso de gestión de vulnerabilidades, un ritmo de divulgación pública y salvaguardas que impidan que la capacidad se reconvierta directamente en automatización de ataques.

\begin{warningbox}{El riesgo de doble uso debe gobernarse desde el principio}
Cuanto más automatizado sea un sistema de descubrimiento de vulnerabilidades, más necesario es convertir en flujo predeterminado el control de acceso, el sandbox de ejecución, la anonimización de resultados y la aprobación de divulgación, en lugar de tratarlos como parches de última hora en el proyecto. De lo contrario, el éxito técnico se convierte directamente en riesgo de cumplimiento normativo.
\end{warningbox}

\begin{importantbox}{Lista mínima de gobernanza de seguridad para la práctica}
\begin{enumerate}
    \item incluir en lista blanca todos los repositorios objetivo y entornos de ejecución;
    \item ejecutar por defecto las entradas de tipo exploit-like solo en un entorno de ejecución aislado;
    \item al detectar una posible vulnerabilidad, verificarla internamente antes de iniciar el proceso de divulgación;
    \item distinguir en el informe entre ``evidencia'' e ``inferencia'', para evitar conclusiones exageradas;
    \item al publicar hacia el exterior, eliminar los detalles sensibles que permitan reproducir directamente la cadena de ataque.
\end{enumerate}
\end{importantbox}

\subsection{Oportunidades de investigación: preguntas que vale la pena abordar a continuación}
Al final, Sutton subraya que este es un ``wide open area''; si se traduce en una agenda de investigación, puede concentrarse en cuatro tipos de problemas: verifiers dinámicos más sólidos, generalización entre repositorios más estable, búsqueda de trayectorias largas con menor costo y generación automática de evidencia más confiable. En conjunto, estos determinan si un Agente de seguridad puede pasar de un éxito puntual a una producción estable y a escala.

\begin{table}[H]
\centering
\begin{tabular}{p{3.0cm}p{4.4cm}p{4.8cm}}
\toprule
\textbf{Dominio del problema} & \textbf{Cuello de botella actual} & \textbf{Posible dirección de avance} \\
\midrule
Verificador dinámico & La señal de activación es inestable y el costo de descartar falsos positivos es alto & Combinar sanitizer + aserciones semánticas + ejecución diferencial \\
Generalización entre proyectos & Las herramientas y los prompts dependen en gran medida de las convenciones de un solo repositorio & Protocolo ACI transferible e interfaz de normalización de tareas \\
Eficiencia de búsqueda & Las trayectorias largas tienen un costo alto y la retroalimentación de los fallos es escasa & Planificación por capas, repetición de experiencia, reutilización de la memoria de fallos \\
Automatización de evidencia & El informe carece de legibilidad y de capacidad de auditoría & Plantillas estructuradas de evidencia de vulnerabilidad y artifacts reproducibles \\
\bottomrule
\end{tabular}
\caption{Mapa de oportunidades de investigación a corto y mediano plazo para Agentes de seguridad}
\end{table}

\subsection{Resumen de la sección}
Trasladar el método de la clase a la práctica de ingeniería depende de dos cosas: primero, toda optimización debe girar en torno a un ciclo cerrado ejecutable; segundo, toda capacidad debe ir acompañada de mecanismos de auditoría y gobernanza. Solo cuando se cumplen a la vez ``poder hacerlo'' y ``poder controlarlo'', un Agente de seguridad puede entrar de manera sostenida en un flujo de producción real.

